% Em coletâneas, use \capitulo[Autores do capítulo]{Título}; em obras autorais, \capitulo{Título}.
\capitulo[Nome do Autor · Nome da Autora]{Título do primeiro capítulo}

\begin{resumo}
Até 150 palavras (coletâneas): objetivo, conteúdo principal e contribuição do capítulo.
\end{resumo}
\palavraschave{termo; termo; termo.}

\section{Seção}
Texto corrido em Source Serif 4, corpo 10,5 pt sobre 14,5 pt, justificado, com recuo de 5 mm na primeira linha. A
hemogasometria\index{Hemogasometria} é usada na rotina de cães\index{Cães} e gatos\index{Gatos}. Para incluir um termo no
índice remissivo, escreva o comando \texttt{\textbackslash index\{Termo\}} logo depois dele; para subentradas, use
\texttt{\textbackslash index\{Acidose!metabólica\}}.

Esta mancha gráfica e esta entrelinha dão de 2.000 a 2.200 caracteres por página no formato 16×23 cm, com boa leitura na tela
e no papel. Títulos em Roboto, na cor da marca. Capítulos sempre abrem em página ímpar, sem cabeçalho.

\begin{citacao}
Citação direta com mais de três linhas, recuada 12 mm, em corpo menor e sem aspas, seguida da indicação da fonte com a página
(Autor, 2024, p. 00).
\end{citacao}

\begin{figure}[htb]
  \centering
  \includegraphics[width=\textwidth]{figura-exemplo}
  \caption{Legenda descritiva da figura, que permita entendê-la sem ler o texto.}
  \fonte{os autores (2026).}
\end{figure}

\begin{table}[htb]
  \caption{Título da tabela, acima dela.}
  \centering\small
  \begin{tabular}{lcc}
    \toprule
    Variável & Grupo A & Grupo B \\
    \midrule
    pH & 7,38 ± 0,02 & 7,21 ± 0,05 \\
    HCO$_3^-$ (mmol/L) & 22 ± 2 & 14 ± 3 \\
    \bottomrule
  \end{tabular}
  \fonte{os autores (2026).}
\end{table}

\section{Outra seção}
A acidose metabólica\index{Acidose!metabólica} é discutida aqui como exemplo de subentrada no índice.

\begin{referencias}
\item SOBRENOME, Nome; SOBRENOME, Nome. Título do artigo. \textbf{Nome da Revista}, v. 0, n. 0, p. 00-00, 2024. DOI: https://doi.org/10.0000/xxxx.
\item SOBRENOME, Nome. \textbf{Título do livro}. 2. ed. Cidade: Editora, 2023.
\end{referencias}
